% 383_Ein_Anker_ohne_v_De_ch.tex
% Johann Pascher, 29. Sept. 2026
% Massen und Planck-Skala ohne v: eine Kette, ein Anker

\providecommand{\BM}{\textbf{[B]}}
\providecommand{\KM}{\textbf{[K]}}
\providecommand{\SM}{\textbf{[S]}}

\chapter*{Massen und Planck-Skala ohne \texorpdfstring{$v$}{v}}
\addcontentsline{toc}{chapter}{Massen und Planck-Skala ohne v}

\begin{abstract}
Die Leptonleiter des Korpus schreibt die Massen als $m_i=r_i\,\xi^{p_i}\,v$ und benutzt
die elektroschwache Skala $v$ als Zwischengröße; die Planck-Skala wird getrennt über
die Formel für $G$ mit dem Umrechnungsfaktor $C_{\text{conv}}$ erreicht. So gelesen
sieht es aus, als brauche die Rechnung zwei Anker, einen für $v$ und einen für $E_P$.
Das ist nicht der Fall. Dok.~149 verbindet $v$ und $E_P$ direkt:
$v=E_P/(f^4\cdot(\pi/2)\cdot10)$, also $v/E_P=\xi^4/(5\pi)=2{,}012\times10^{-17}$,
$-0{,}23\,\%$ gegen den gemessenen Wert \KM. Setzt man das in die Leiter ein, fällt
$v$ heraus: $m_i=\frac{r_i}{5\pi}\,\xi^{p_i+4}\,E_P$ \BM. Damit hängen alle drei
Leptonmassen, $E_P$, $G$, $\ell_P$ und $L_0$ an einer einzigen Kette, und ein einziger
Messwert als Anker legt alle Zahlen fest. Mit $m_e$ als Anker folgen
$E_P=1{,}237\times10^{19}$\,GeV ($+1{,}34\,\%$), $G=6{,}50\times10^{-11}$
($-2{,}6\,\%$), $\ell_P=1{,}595\times10^{-35}$\,m und $L_0=2{,}127\times10^{-39}$\,m
\KM. Die Abweichungen sind die bekannten bare-Reste der Leiter, kein zweiter Anker.
\end{abstract}

\noindent\textbf{Zur Bezeichnung:} $f$ ist hier die Grundwindungszahl $f = 1/\xi = 7500$, keine Frequenz. In natürlichen Einheiten ist $f$ nur eine andere Lesart der Dualität $T\cdot m = 1$: $f = t_P/t_0 = \Lambda_{\text{T0}}/E_P$ mit $t_0 = \xi\,t_P$ und $\Lambda_{\text{T0}} = E_P/\xi$, und aus $\xi E_0^2 = 1$ folgt $f = E_0^2$. Als dimensionsloses Verhältnis lässt sich $f$ nicht begründet auf~1 setzen (Dok.~A135).

% ============================================================
\section*{Statusmarkierungen}
% ============================================================
\begin{center}
\begin{tabular}{@{}lp{10cm}@{}}
\textbf{[B]} & Algebraisch bewiesen (Umformung ohne Zusatzannahme). \\[3pt]
\textbf{[K]} & Numerisch bestätigt (Prüfskript). \\[3pt]
\textbf{[S]} & Spekulativ: Hinweis oder Vermutung, im Korpus noch nicht hergeleitet oder geprüft. \\
\end{tabular}
\end{center}
Messwerte (CODATA 2022, PDG) dienen als Anker bzw.\ als Vergleichswerte. $\hbar$, $c$
und $e$ sind im SI exakt festgelegt; sie bestimmen nur die Einheiten und sind kein
zusätzlicher Anker.

% ============================================================
\section{Ausgangslage}
% ============================================================
Im Korpus stehen zwei Wege nebeneinander:
\begin{enumerate}
\item \textbf{Leptonleiter} (Dok.~006, 352): $m_i=r_i\,\xi^{p_i}\,v$ mit
  $r=(4/3,\;16/5,\;25/9)$ und $p=(3/2,\;1,\;2/3)$ für $e,\mu,\tau$. Die
  Massenverhältnisse $m_\mu/m_e=\tfrac{12}{5}\xi^{-1/2}=207{,}85$ und
  $m_\tau/m_\mu=\tfrac{125}{144}\xi^{-1/3}=16{,}99$ enthalten $v$ nicht.
\item \textbf{Gravitation} (Dok.~012, 013, 180):
  $G=\xi^2/(4m_e)\cdot C_{\text{conv}}\cdot K_{\text{frak}}$, daraus $E_P$ und $\ell_P$.
\end{enumerate}
Wer $v$ aus $m_e$ gewinnt und $E_P$ aus der $G$-Formel, verwendet zwei Messwerte:
$m_e$ und, über $C_{\text{conv}}$, den Messwert von $G$. So rechnet Dok.~375. Eine
Kette, die aus $\xi$ allein alle Größen herleiten soll, darf aber nur einen Anker
haben; alle übrigen Zahlen müssen folgen. Das Bindeglied liefert Dok.~149.

% ============================================================
\section{Die Beziehung zwischen \texorpdfstring{$v$}{v} und \texorpdfstring{$E_P$}{E\_P}}
% ============================================================
Dok.~149 leitet die elektroschwache Skala aus der Planck-Energie ab: Verdünnung über
die $f^4$ Zellen des vierdimensionalen Gitters, Projektion auf den dreidimensionalen
Halbraum ($\pi/2$) und ein Faktor $1/10$ für den Übergang zur elektroschwachen Skala:
\begin{equation}
  v=\frac{E_P}{f^4\cdot(\pi/2)\cdot10}
  \qquad\Longleftrightarrow\qquad
  \frac{v}{E_P}=\frac{\xi^4}{5\pi}\quad\BM
\end{equation}
mit $f=1/\xi$. Die Umformung ist exakt. Der Faktor~$10$ ist weder hergeleitet noch
motiviert, sondern an den Messwert angepasst: Im Arbeitsstand vom Februar 2026 ergab
$m_P/(f^4\cdot\pi/2)$ den Wert $2467$~GeV statt $246$~GeV, und die Lücke wurde mit
$1/10$ geschlossen (Dok.~149). Setzt man statt des gemessenen $v$ das nackte
$v=248{,}93$~GeV der Leiter ein, verlangt die Formel $10\,K_{\text{frak}}\approx\pi^2$
statt~10; der Faktor~10 wäre dann eine Vermischung von nacktem Leiterwert und
SI-Vergleichswert. Das ist eine Setzung und offen \SM; die Varianten vergleicht
Dok.~387.

\textbf{Zahlenwert.} $v/E_P=\xi^4/(5\pi)=2{,}0120\times10^{-17}$. Gemessen ist
$v/E_P=246{,}22\,\text{GeV}/1{,}22089\times10^{19}\,\text{GeV}=2{,}0167\times10^{-17}$;
die Abweichung beträgt $-0{,}23\,\%$ \KM. Weil $v/E_P$ dimensionslos ist, braucht
diese Zahl keinen Anker.

\textbf{Zum Wert in Dok.~149.} Dort wird mit $f=7491{,}91$ gerechnet und
$v=246{,}71$\,GeV angegeben. Mit $f=1/\xi=7500$ ergibt dieselbe Formel
$v=245{,}65$\,GeV \KM. Beide Werte sind richtig gerechnet; sie unterscheiden sich im
eingesetzten $f$.

% ============================================================
\section{Die Massen ohne \texorpdfstring{$v$}{v}}
% ============================================================
Einsetzen von $v=\xi^4E_P/(5\pi)$ in die Leiter ergibt
\begin{equation}
  m_i=\frac{r_i}{5\pi}\;\xi^{\,p_i+4}\;E_P\quad\BM
\end{equation}
Mit dem Messwert $E_P=1{,}22089\times10^{22}$\,MeV als Vergleich:

\begin{center}
\begin{tabular}{@{}lllrr@{}}
\toprule
\textbf{Lepton} & \textbf{Formel} & \textbf{Koeffizient} & \textbf{Wert} & \textbf{gegen Messung}\\
\midrule
$e$    & $\frac{4}{15\pi}\,\xi^{11/2}E_P$ & $0{,}08488$ & $0{,}5043$\,MeV  & $-1{,}32\,\%$\\[3pt]
$\mu$  & $\frac{16}{25\pi}\,\xi^{5}E_P$   & $0{,}20372$ & $104{,}81$\,MeV  & $-0{,}80\,\%$\\[3pt]
$\tau$ & $\frac{5}{9\pi}\,\xi^{14/3}E_P$ & $0{,}17684$ & $1780{,}9$\,MeV  & $+0{,}23\,\%$\\
\bottomrule
\end{tabular}
\end{center}
\KM\ Die Verhältnisse der Massen sind dieselben wie in der Leiter, denn $E_P$ und
$5\pi$ kürzen sich heraus. Die Abweichungen setzen sich zusammen aus dem bare-Rest der
Leiter gegen $v$ (Dok.~352) und den $-0{,}23\,\%$ der Beziehung $v/E_P$.

% ============================================================
\section{Eine Kette, ein Anker}
% ============================================================
Die Kette lautet:
\begin{equation}
  \xi\;\longrightarrow\;\frac{m_i}{E_P}=\frac{r_i}{5\pi}\,\xi^{p_i+4}
  \;\longrightarrow\;E_P\;\longrightarrow\;
  G=\frac{\hbar c^5}{E_P^2},\quad
  \ell_P=\frac{\hbar c}{E_P},\quad
  L_0=\xi\,\ell_P .
\end{equation}
Bis auf die Einheiten ist alles durch $\xi$ bestimmt. Ein einziger Messwert legt die
Zahlen fest. Mit $m_e=0{,}51099895$\,MeV als Anker:

\begin{center}
\begin{tabular}{@{}lrr@{}}
\toprule
\textbf{Größe} & \textbf{aus Anker $m_e$} & \textbf{gegen Messung}\\
\midrule
$E_P$    & $1{,}2372\times10^{19}$\,GeV & $+1{,}34\,\%$\\
$G$      & $6{,}4995\times10^{-11}$\,m$^3$\,kg$^{-1}$\,s$^{-2}$ & $-2{,}6\,\%$\\
$\ell_P$ & $1{,}5950\times10^{-35}$\,m & $-1{,}32\,\%$\\
$L_0=\xi\ell_P$ & $2{,}1266\times10^{-39}$\,m & $-1{,}32\,\%$\\
\bottomrule
\end{tabular}
\end{center}
\KM\ Da $G\propto E_P^{-2}$, ist die relative Abweichung von $G$ doppelt so groß wie
die von $E_P$.

\textbf{Wahl des Ankers.} Die Kette kann an jeder Stelle beginnen; welche Größe man
als Anker nimmt, ist eine Frage der Genauigkeit, nicht der Herleitung. Die
Abweichungen der übrigen Größen hängen dann vom bare-Rest des gewählten Einstiegs ab:

\begin{center}
\begin{tabular}{@{}lrrr@{}}
\toprule
\textbf{Anker} & $E_P$ & $G$ & $\ell_P$, $L_0$\\
\midrule
$m_e$    & $+1{,}34\,\%$ & $-2{,}62\,\%$ & $-1{,}32\,\%$\\
$m_\mu$  & $+0{,}81\,\%$ & $-1{,}60\,\%$ & $-0{,}80\,\%$\\
$m_\tau$ & $-0{,}23\,\%$ & $+0{,}45\,\%$ & $+0{,}23\,\%$\\
$v$      & $+0{,}23\,\%$ & $-0{,}46\,\%$ & $-0{,}23\,\%$\\
\bottomrule
\end{tabular}
\end{center}
\KM\ In jedem Fall wird genau ein Messwert verwendet. Die Streuung zwischen den Zeilen
ist der bare-Rest der Leiter; sie verschwindet in dem Maß, in dem Dok.~352 diesen Rest
auflöst.

% ============================================================
\section{Verhältnis zur \texorpdfstring{$G$}{G}-Formel mit \texorpdfstring{$C_{\text{conv}}$}{C\_conv}}
% ============================================================
Die Formel $G=\xi^2/(4m_e)\cdot C_{\text{conv}}\cdot K_{\text{frak}}$ aus Dok.~012
ergibt $6{,}6745\times10^{-11}$, also $+0{,}003\,\%$. Das ist kein Widerspruch zu den
$-2{,}6\,\%$ oben: $C_{\text{conv}}=7{,}783\times10^{-3}$ ist so gewählt, dass die
SI-Umrechnung den gemessenen Wert von $G$ trifft. In dieser Form gehen $m_e$ und $G$
als Anker ein. Die Ein-Anker-Kette mit $m_e$ entspricht
$C_{\text{conv}}=7{,}58\times10^{-3}$ \KM; der Unterschied von $2{,}6\,\%$ ist genau
der Rest der Kette. $C_{\text{conv}}$ bleibt die SI-Umrechnung von $G$; die Kette zeigt
nur, dass man ihn nicht als zweiten Messwert braucht, wenn man den Rest in Kauf nimmt.

Dok.~375 rechnet $v/E_P$ auf diesem Weg ($v$ aus $m_e$, $E_P$ über $C_{\text{conv}}$)
und erhält $2{,}0389\times10^{-17}$. Der direkte Wert $\xi^4/(5\pi)=2{,}0120\times10^{-17}$
liegt um den Faktor $1{,}0134$ darunter; das ist genau der Rest $-1{,}32\,\%$ der
Elektronmasse in der Kette \KM. Beide Wege sind damit konsistent.

% ============================================================
\section{Bezüge in anderen Dokumenten}
% ============================================================
\begin{itemize}
\item \textbf{Dok.~149:} An der Formel für $v$ steht, dass sie mit $f=7500$
  $245{,}65$\,GeV ergibt und $v/E_P=\xi^4/(5\pi)$ bedeutet.
\item \textbf{Dok.~180:} Bei $\ell_P$ steht, dass der Wert $1{,}616\times10^{-35}$\,m
  der Messwert als Anker ist und die Ein-Anker-Kette $1{,}595\times10^{-35}$\,m ergibt.
\item \textbf{Dok.~375:} Dort steht, dass $v/E_P$ direkt als $\xi^4/(5\pi)$ folgt, ohne
  $C_{\text{conv}}$ und ohne Anker.
\end{itemize}

% ============================================================
\section{Ergebnis}
% ============================================================
$v$ ist in der Massenrechnung eine Zwischengröße und lässt sich auflösen:
$m_i=\frac{r_i}{5\pi}\xi^{p_i+4}E_P$ \BM. Damit hängen die Leptonmassen, $E_P$, $G$,
$\ell_P$ und $L_0$ an einer Kette, die bis auf die Einheiten durch $\xi$ bestimmt ist.
Ein einziger Messwert genügt als Anker; welcher, ist frei. Die Abweichungen von
$0{,}2$ bis $1{,}3\,\%$ in $E_P$ sind die bare-Reste der Leiter und der Beziehung
$v/E_P$ ($-0{,}23\,\%$), keine Folge eines fehlenden Ankers. Der Faktor $10$ in dieser
Beziehung ist die offene Stelle der Kette \SM.

\medskip\noindent Prüfskript: \texttt{2/python/Dok383\_Skripte/pruef\_383\_ein\_anker.py}, 20/20.

\begin{thebibliography}{9}
\bibitem{CODATA2022} E.~Tiesinga et al., \emph{CODATA Recommended Values of the
  Fundamental Physical Constants: 2022}, NIST (2024).
\bibitem{PDG2024} Particle Data Group (S.~Navas et al.), \emph{Review of Particle
  Physics}, Phys.\ Rev.\ D 110 (2024) 030001.
\end{thebibliography}
